\documentclass[11pt]{article}

\usepackage[margin=1in]{geometry}
\usepackage[T1]{fontenc}
\usepackage[utf8]{inputenc}
\usepackage{lmodern}
\usepackage{microtype}
\usepackage[round]{natbib}
\usepackage{hyperref}
\usepackage{enumitem}

\hypersetup{
  colorlinks=true,
  linkcolor=blue,
  citecolor=blue,
  urlcolor=blue
}

\setlist[itemize]{leftmargin=1.5em}
\setlist[enumerate]{leftmargin=1.5em}

\title{Death of the Author, Birth of the Scriptor:\\ \large AI in Academia, Incentives, and Post-Paper Research Outputs\protect\footnotemark[1]\protect\footnotemark[2]}
\author{}
\date{May 2026\protect\footnotemark[3]}

\begin{document}
\maketitle
\footnotetext[3]{Draft dated 13 May 2026.}

\begin{abstract}
Cesar Hidalgo warns that an AI tsunami is approaching science. The warning matters because AI lowers the execution costs of traditional research tasks: search, coding, modelling, drafting, review, visualisation, and revision. In weak-link terms, AI moves the constraint on research productivity towards specifying, auditing, validating, and warranting the workflow through which claims are produced. This challenges the academic paper as the authoritative way of knowing research. The paper should become the narrative layer of a wider evidence system: a modular constellation of claims, sources, code chains, prompts, checks, exclusions, revisions, uncertainty, and public outputs. Quantity-output governance, patched with ad hoc data, code, and AI-use statements, is poorly suited to rigorous research under generative conditions. Where a prompter uses AI to produce papers, a scriptor recognises that the workflow is the output. \texttt{scriptorship-ledger} shows this shift through a claim ledger, transparency log, source-bias report, warrant graph, validation scripts, and workflow profile, paired with a limited v1 Codex/LaTeX pipeline that researches, validates, tracks accepted edits, and makes the writing workflow inspectable.
\vspace{5.5cm}
\end{abstract}
\footnotetext[1]{A public workflow repository archives the LaTeX project, claim ledger, transparency log, validation scripts, Overleaf/Git workflow, and prototype reviewer tooling. Codex helped formalise an existing paper-pipeline approach; the initiating prompt is retained as part of the workflow record.}
\footnotetext[2]{The title invokes \citet{barthes1977death} to mark a move away from treating the named author as the sovereign origin that fixes meaning. In AI-mediated research, authorial status remains a responsibility category, while claim-level trust has to move through the workflow. For recent discussions of Barthes, authorship, and AI-generated work, see \citet{masters2025academicAuthor} and \citet{puscasu2024creatingAI}.}
\setcounter{footnote}{3}

\pagebreak
\input{sections/01_introduction}
\input{sections/02_output_based_trust}
\input{sections/03_falling_execution_costs}
\input{sections/04_birth_of_the_scriptor}
\input{sections/05_modularisation_research_knowledge}
\input{sections/06_script_protocol}
\input{sections/07_claim_level_warranting}
\input{sections/08_verification}
\input{sections/09_incentives}
\input{sections/10_authorship_voice_workflow_transparency}
\input{sections/11_implementation}
\input{sections/12_conclusion}

\pagebreak
\bibliographystyle{plainnat}
\bibliography{references}

\end{document}
